\documentclass[10pt]{article}
\usepackage[margin=20mm]{geometry}
\usepackage{amsmath,amssymb,mathtools,fontspec,unicode-math,xurl,hyperref}
\IfFontExistsTF{DejaVu Serif}{\setmainfont{DejaVu Serif}}{\setmainfont{Latin Modern Roman}}
\setmathfont{latinmodern-math.otf}
\setlength{\parindent}{0pt}\setlength{\parskip}{6pt}\sloppy
\title{Initial mathematical review: M25-0157
On the Balmer spectrum of the Morel–Voevodsky category}
\author{}
\date{26 October 2026}
\begin{document}
\maketitle
Authors: Peng Du and Alexander Vishik. Publication identity: Duke Mathematical Journal 174 (2025), no. 6, 1013–1044; DOI 10.1215/00127094-2024-0051 (https://doi.org/10.1215/00127094-2024-0051).\par
Reviewed version: arXiv:2309.09077v3, 28 July 2024, 28 PDF pages. All page references below are to that manuscript, with its printed pagination.\par
Initial assessment: Major. Two printed derivations supporting the main theorem use false auxiliary assertions: the infinite-height bridge to conservativity, and annihilation of a local projector in the field-equivalence argument. Both required intermediate conclusions admit the repairs proved below; no counterexample to Theorem 7.6 is established. A separate finite-generation assertion also fails at infinite height, outside its actual downstream use. Severity is assigned to the demonstrated proof dependencies, separately from survival of the repaired conclusions. This is an initial review under frozen contract 0.7.1; the scientific decision remains `pending`, and `human\_reviews=[]`. It is not an independent acceptance decision.\par
\section*{Scope and central argument}
The paper constructs prime tensor ideals \(\mathfrak a_{(p,m),E}\) of the compact Morel–Voevodsky category over a characteristic-zero field. Here \(p\) is prime, \(1\le m\le\infty\), and \(E/k\) is an arbitrary field extension, considered up to the stated \(K(p,m)\)-equivalence. Theorem 7.6 asserts primality, the exact coincidence criterion for these ideals, and compatibility at \(m=\infty\) with isotropic motivic points. The use of a flexible closure, compactness, and the infinite-height convention are substantive hypotheses, not cosmetic notation.\par
The entire assigned text was read, and the following chain was examined: anisotropy and its coefficient-ideal criterion (Section 2); the two localizations and the compact-object description (Section 3); the bounded weight structure and numerical heart (Section 4); invariant annihilator radicals (Section 5); semisimplicity and the tensor weight complex (Section 6); the finite-height annihilator argument, primality, and the infinite-height endpoint (Section 7); and local categories, conservativity, separation of fields and heights, and examples (Section 8). They were collaborators within this initial review, not independent adjudicators.\par
The exact coefficient convention matters: \[\mathbb L_{(p)}=BP[x_i:i\ne p^r-1],\qquad
 P(m)=BP/(p,v_1,\ldots,v_{m-1}).\] At infinite height \(P(\infty)=\mathbb F_p\), but \(P(\infty)[x]=\mathbb F_p[x]\) still has the non-\(p\)-typical polynomial variables. The target of the Morava coefficient change is \(K(p,\infty)=\mathbb F_p\).\par
\section*{Problem 1: finite generation fails at infinite height}
Location. Proposition 5.2, p. 12, and its finite-presentation argument on p. 13; compare Proposition 4.3, p. 11, and the separate infinite-height case in Proposition 7.4, p. 17.\par
Claim and decisive test. Proposition 5.2 states that the radical of the coefficient annihilator of a compact object is finitely generated. It does not restrict \(m\) to finite height. Take a flexible characteristic-zero field, for example \(k=\mathbb C(t_1,t_2,\ldots)\), and take the compact localized unit in \(\widehat{\mathrm{SH}}_{(p,\infty)}(k/k)\). This is the first localization, before the further quotient in Definition 3.1. Its \(\mathrm{MGL}\)-motive is the unit in the numerical heart of Proposition 4.3. The graded coefficient endomorphisms of that unit are \(\mathbb F_p[x]\). Numerical equivalence introduces no additional relation on the point: pairing a nonzero coefficient with \(1\) detects it. Consequently, after \(p\)-localization, \[J=\operatorname{Ann}_{\mathbb L_{(p)}}(\widehat M^{\mathrm{MGL}}(\mathbf1))
   =(p,v_1,v_2,\ldots)\mathbb L_{(p)} .\] This ideal is radical, since the quotient is the polynomial domain \(\mathbb F_p[x]\).\par
It is not finitely generated. Any finite family of elements of \(J\) involves only finitely many of the variables \(v_i\). For a sufficiently large \(N\), that family belongs to \((p,v_1,\ldots,v_N)\). Its generated ideal therefore cannot contain \(v_{N+1}\): the latter survives in the polynomial quotient by \((p,v_1,\ldots,v_N)\). This is an algebraic counterexample under the paper’s own flexible-field and compactness hypotheses. Finite generation before localization would imply finite generation after localization, so working \(p\)-locally does not evade the claim.\par
Failed proof step. For the one-term weight complex of this unit, the endomorphism module is \(\mathbb F_p[x]\) as an \(\mathbb L_{(p)}\)-module. It has infinitely many independent defining relations \(p,v_1,v_2,\ldots\). The finite-presentation assertion used on p. 13 is therefore unavailable at \(m=\infty\).\par
Impact, counterevidence, and repair boundary. Restrict Proposition 5.2 and the finite-presentation argument to \(m<\infty\), with the flexible-field hypothesis used in the proof made explicit. The demonstrated counterexample does not affect finite height. More importantly, Proposition 7.4 explicitly separates \(m=\infty\) and uses conservativity there, rather than this finite-generation argument. The downstream main-theorem damage is thus not established. The defect is assessed as Minor, based on this actual dependency boundary. The unit counterexample itself is fully verified; a uniform theorem over nonflexible fields is not supplied by this review.\par
There is also a harmless logical convention to state in Proposition 5.3: if the \(\mathrm{MGL}\)-motive is zero, its annihilator is the whole ring. The proper-prime description applies to the nonzero case; the zero case in Proposition 7.4 is immediate. This qualification is recorded without promoting it to a separate Problem.\par
\section*{Problem 2: the infinite-height functor is not an equivalence}
Location. The coefficient-change functor displayed immediately before Proposition 7.1, p. 15, and the first sentence of its proof, p. 16. Its required downstream use is the last sentence of Proposition 7.4, p. 17.\par
Decisive evidence. At \(m=\infty\) this functor is induced by \[R=\mathbb F_p[x]\longrightarrow \mathbb F_p,\qquad x_i\longmapsto0.\] It is not faithful, even on the heart, and hence cannot be an equivalence. For instance, take \(p=2\). The allowed variable \(x_2\) is a nonzero coefficient morphism \(\mathbf1(2)[4]\to\mathbf1\) in numerical \(R\)-motives. It survives numerical equivalence over the point, but maps to zero after coefficient change. Viewed as a morphism between one-term complexes, it gives the same failure in the bounded homotopy category. Thus the asserted equivalence cannot be justified by passing to complexes.\par
The needed weaker assertion and its repair. The argument needs conservativity, not faithfulness. Here is a bounded proof of that weaker assertion. Numerical correspondences for the polynomial extension of \(CH/p\) are numerical \(CH/p\) correspondences tensored with \(R\): numerical triviality is tested coefficient by coefficient by the degree pairing. In the presentation inherited from \(\mathbb L_{(p)}=BP[x]\), an orientation change gives the same correspondence category; its Todd factors are invertible and do not remove the polynomial variables.\par
Let \(V\) be a fixed finite direct sum of twisted smooth projective motives, and put \(J_x=(x)\). The kernel of \[\operatorname{End}(V)\longrightarrow\operatorname{End}(V\otimes_R R/J_x)\] is nilpotent. Each factor in a product from this kernel contributes positive coefficient degree. In a fixed source–target correspondence group, the compensating Chow codimension is bounded by the dimensions and the fixed twists. A sufficiently long product therefore vanishes. Reduction on Hom groups is surjective, and these properties pass to idempotent summands.\par
Now let a bounded complex \(C\) become zero in the target homotopy category. Choose a contracting homotopy there and lift its components to a degree \(-1\) map \(h\) on \(C\). The chain endomorphism \(q=dh+hd\) reduces to \(\operatorname{id}\). The error \(\delta=\operatorname{id}-q\) lies in the preceding nilpotent kernel on the finite direct sum of the terms of \(C\). Hence \[q^{-1}=\operatorname{id}+\delta+\cdots+\delta^{N-1}\] exists for some \(N\), commutes with \(d\), and \(hq^{-1}\) contracts \(C\). This proves conservativity. Combining it with conservativity of the bounded weight complex yields exactly the infinite-height implication needed in Proposition 7.4.\par
Impact and classification. The nonfaithful Tate map alone refutes only the stronger equivalence assertion. The core proof problem is more specific: bounded-weight conservativity applies first to the weight complex over \(\mathbb F_p[x]\), whereas Proposition 7.4 needs conservativity after further reduction to \(\mathbb F_p\). The false equivalence on p. 16 is the manuscript’s stated bridge across this loss of coefficients; the final sentence on p. 17 invokes the resulting conservativity without a separate justification. That route feeds Proposition 7.5 and Theorem 7.6(1),(3) at infinite height. The contraction argument above supplies the essential missing bridge. Under contract 0.7.1, the demonstrated invalid derivation directly supporting this part of the main contribution is Major. Its repairability does not lower that assessment. The repaired conservativity is true, and no counterexample to Proposition 7.4 or Theorem 7.6 is claimed.\par
\section*{Problem 3: the bare local projector need not be killed by \(p\)}
Location. Proof of Proposition 8.5, p. 24, after the first correspondence diagram: the assertion that the identity of \(\widetilde{\mathcal X}_{Q}\) is killed by \(p\). Here \(Q\) comprises smooth projective varieties whose function fields strictly dominate the chosen field under the \(K(p,m)\)-order. The actual subsequent maps also contain \(M^{\mathrm{MGL}}(R)\) or \(M^{\mathrm{MGL}}(P)\) as tensor factors.\par
Counterexample to the stated intermediate assertion. Take \(p=2\), finite \(m\ge1\), \[k=\mathbb Q(t_1,t_2,\ldots),\qquad E=F=k(\sqrt2),\qquad
 P=R=\operatorname{Spec}E .\] This satisfies the flexible-base and finitely generated equivalent-extension conditions in the proof. For every integral component \(Q_j\) of the defining union, \(k(Q_j)\ge_{(2,m)}E\) forces \(E\subset k(Q_j)\). Indeed, the quadratic finite scheme \(\operatorname{Spec}E\) over \(k(Q_j)\) is \(K(2,m)\)-isotropic only if it splits. Normality extends this field inclusion to an \(E\)-structure on \(Q_j\).\par
Strict domination means that \((Q_j)_E\) is \(K(2,m)\)-anisotropic. Every closed point of \(Q_j\), as an \(E\)-variety, consequently has even degree: an odd-degree point would have nonzero pushforward in characteristic two. Every closed-point degree over \(k\) is therefore divisible by four. Conversely, the defining union includes \(\operatorname{Spec}E(\sqrt{t_1})\), which has degree four over \(k\) and strictly dominates \(E\). Its index ideal over \(k\) is thus exactly \(4\mathbb Z\).\par
Apply the monoidal change from \(\mathrm{MGL}\)-modules to integral motives. For the complementary Čech projector, the degree-zero homology sequence gives \[H_{0,0}\bigl(\widetilde M(\check C(Q));\mathbb Z\bigr)
 \cong \mathbb Z/\operatorname{im}\bigl(H_{0,0}(M(\check C(Q)))\to\mathbb Z\bigr)
 =\mathbb Z/4 .\] The image is precisely the degree ideal; negative degree-zero-weight homology of the effective Čech complex vanishes. Its unit class is not killed by \(2\). Thus \(2\operatorname{id}\) cannot vanish on the bare projector in \(\mathrm{MGL}\)-modules either: vanishing would survive change of coefficients.\par
Required tensor factors and repair. The needed assertions are \[p\,\operatorname{id}_{M^{\mathrm{MGL}}(R)\otimes\widetilde{\mathcal X}_Q}=0,
 \qquad
 p\,\operatorname{id}_{M^{\mathrm{MGL}}(P)\otimes\widetilde{\mathcal X}_Q}=0 .\] Choose the degree-\(p\) extension \(L/k\) in the proof, remaining anisotropic over both function fields. The finite covers \(R_L\to R\) and \(P_L\to P\) have degree \(p\), and their function fields strictly dominate the respective fields. Their motives are killed by the complementary projector. Pullback followed by pushforward along either cover is multiplication by \(p\), proving the displayed tensor-factor versions. On these objects multiplication by an integer prime to \(p\) is invertible. This is sufficient to remove the integer \(l\) in the two factorization maps on p. 24.\par
Impact and repair boundary. The bare-projector statement is explicitly used to remove \(l\) and obtain the two \(v_m\)-factorizations. These factorizations establish equality of kernels in Proposition 8.5, which supplies the field-equivalence part of Theorem 7.6(2). Thus the false assertion is an operative justification in a core derivation, not an unused stronger aside. This meets the Major criterion in contract 0.7.1. Both maps in fact contain the required \(P\) or \(R\) factor, so the replacement transfer argument proves their needed form and restores this step. The correction is local in extent, but its brevity does not change the severity of the printed derivation. The counterexample is not a counterexample to equality of the two kernels or to Theorem 7.6.\par
\section*{Checks supporting the unaffected argument}
At finite height, the compactness in Proposition 3.3 is essential: a compact object in the localizing ideal generated by the compact cone of \(v_m\) lies in its thick closure, and hence is killed by a finite power of \(v_m\). Conversely a nilpotent action makes the object a retract of a cone construction. The review found no reversal of this implication.\par
The finite-height proof of Proposition 7.1 was checked separately from the defective endpoint sentence. Localization in \(v_m\) kills the higher Landweber factors and leaves finite free modules. Lemma 7.2 is the determinant-trick annihilator argument; in Lemma 7.3 an annihilator acting trivially on Tor and as a unit after coefficient change forces Tor to vanish. In the spectral sequence on p. 17, the induction starts at the lowest nonzero homology degree. Incoming differentials to the edge term originate in already vanishing lower rows, and outgoing differentials have negative first index. This makes the stated induction valid.\par
The primality argument in Proposition 6.6 uses semisimplicity of numerical Morava motives and the absence of tensor zero divisors. The numerical pairing detects products over the coefficient domain; bounded complexes split into shifted cohomology objects. The strong tensor weight complex, rather than an arbitrary weak complex, is the appropriate bridge. The relevant statements of Aoki and Sosnilo were checked, and the bounded-weight kernel-nilpotence input of Bondarko was inspected. Landweber’s original paper was not obtained as a readable full text; that classical input remains an explicitly identified external reliance.\par
The quantifiers in the order on arbitrary field extensions were checked against the finitely generated reductions. The local-to-isotropic conservativity statement in Proposition 8.3 is used on the image of compact objects from the base field, not asserted for every object without that restriction. The cited torsion-space results supply the test objects separating heights; their invariance under field extension is essential. Norm-variety examples give the stated direction of height separation. A reversed isotropy implication in the prose before Example 8.9 is not used to invalidate the correctly directed conclusion and is not counted as a separate scientific Problem.\par
These checks are targeted mathematical audits. They do not amount to reconstructing all of the cited literature, a formal verification of motivic category foundations, or proof of every ancillary example. No unsupported defect is inferred from absent code.\par
\section*{Disposition and verification state}
Item | Assessment | Repair verification | Author adoption | Remaining work\par
Problem 1 | Minor; explicit infinite-height counterexample | Restriction to finite height matches the actual branch in 7.4 | Not observed | Revise 5.2’s scope and handle the zero-motive case explicitly\par
Problem 2 | Major; invalid core infinity derivation | Conservativity proved by nilpotent-kernel lifting of a contraction | Not observed | Replace equivalence and include the needed proof\par
Problem 3 | Major; invalid field-equivalence step | Transfer proves annihilation after the required tensor factors | Not observed | Amend the assertion on p. 24\par
Main theorem | No counterexample established in the checked scope | Listed local repairs preserve the audited dependency chain | Not observed | Scientific decision pending; independent human review remains\par
\section*{Supporting checks}
The source manifest contains three actual input files: the supplied PDF (331703 bytes), the arXiv source archive (27334 bytes), and the extracted author TeX (92753 bytes). All were hashed before review and checked again at delivery. The PDF SHA256 is\par
The source TeX SHA256 is\par
The source manifest SHA256 is\par
The arXiv record and the author’s publication index were checked for identity and available versions. The Nottingham repository advertises an accepted manuscript, but its page was inaccessible; the Project Euclid download resolved to an access-check page. Neither was treated as an examined journal body. No assertion is made that these defects persist in the final journal version.\par
Primary dependencies actually inspected in bounded parts include A. Vishik, Isotropic and numerical equivalence for Chow groups and Morava K-theories, arXiv:2307.15148v2; On torsion spaces, arXiv:2308.14318v1; Isotropic motives, arXiv:1904.09263v3; M. Bondarko, arXiv:0704.4003v9; K. Aoki, arXiv:1904.01384v1; and V. Sosnilo, Theorem of the heart in negative K-theory for weight structures, Documenta Mathematica 24 (2019). The saved notes distinguish manuscript and published numbering.\par
The report is real TeX compiled by the existing Tectonic 0.17.0 executable in untrusted mode. The final compiler command, binary hash, logs, exact TeX before/after hashes, PDF hash, retained dependency inventory, and rendered-page hashes are in the accompanying verification records. No author code or shell escape was executed.\par
Bounded native assistance was provided by `/root/meanfield\_check` (weights and annihilator radicals), `/root/appendix\_check` (semisimplicity, tensor weight complex, and the infinite-height repair), and `/root/resistance\_check` (local categories, point separation, and the projector correction). The primary unit evaluated and integrated their findings. Controller-observed UI information is recorded as such; this worker did not independently verify a backend model or UI selection.\par
Delivery status is distinct from scientific acceptance. A complete compiled and inspected initial-review package does not resolve the pending scientific decision or establish author adoption of any repair.\par
\end{document}
